\documentclass[a4paper,12pt]{article}
\usepackage{docsTemplate}
\usepackage{fontawesome5}

\setTitle{Manuale Utente}
\setAuthors{Sara Gioia Fichera \\ & Andrea Masiero \\ & Filippo Compagno}
\setVerificators{Nenad Radulovic \\ & Francesco Balestro \\ & Bianca Zaghetto}
\setApprovation{}
\setVersion{0.5}
\setType{Esterno}
\setDestination{Team Three Technologies \\ & Prof. Tullio Vardanega \\ & Prof. Riccardo Cardin \\ & Sanmarco Informatica SpA}

\begin{document}
    \include{docTitlepage}

    \section*{Registro delle modifiche} {
        \begin{table}[h!]
            \rowcolors{2}{lightgray}{white}
                \begin{tabularx}{\textwidth}{|l|l|l|l|L|L|}
                \hline
                \textbf{Vers.} & \textbf{Data} & \textbf{Autore} & \textbf{Ruolo} & \textbf{Verifica} & \textbf{Oggetto} \\
                \hline
                \textbf{0.5} & 24/04/26 & Filippo Compagno & Programmatore & Bianca Zaghetto & Aggiornata sezione Guida all'Utilizzo\\
                \hline
                \textbf{0.4} & 14/04/26 & Mattia Oliva Medin & Responsabile & Francesco Balestro & Aggiornata sezione Guida all'Utilizzo\\
                \hline
                \textbf{0.3} & 13/04/26 & Andrea Masiero & Amministratore & Francesco Balestro & Stesura sezione Installazione\\
                \hline
                \textbf{0.2} & 11/04/26 & Sara Gioia Fichera & Responsabile & Nenad Radulovic & Stesura sezione Guida all'Utilizzo\\
                \hline
                \textbf{0.1} & 11/04/26 & Sara Gioia Fichera & Responsabile & Nenad Radulovic & Struttura generale del documento\\
                \hline
            \end{tabularx}
        \end{table}
    }
    \newpage
    
    \tableofcontents
    \newpage

    \listoffigures
    \newpage
    
    \listoftables
    \newpage

    \section{Introduzione} {
        \subsection{Scopo del documento} {        
            Il presente documento ha lo scopo di fornire all'utente tutte le istruzioni necessarie per la corretta installazione e il normale utilizzo dell'applicativo desktop DIPReader.
            Al suo interno è presente una guida per la consultazione dei pacchetti informativi (DIP).
        }

        \subsection{Glossario}
            I termini tecnici o specifici del dominio che potrebbero risultare ambigui sono definiti nel documento \textit{Glossario}, consultabile tra i riferimenti in apertura del documento. All'interno del testo, tali termini sono 
            riconoscibili dalla notazione:
                \begin{center}
                    parola\textsubscript{G}
                \end{center}
        
        \subsection{Riferimenti} {
            \subsubsection{Riferimenti normativi} {
            \begin{itemize}
                \item \textit{Norme di Progetto v1.0}
                \begin{itemize}
                    \item \url{https://team-three-technologies.github.io/Documenti/1%20-%20RTB/Documenti%20interni/Norme_di_Progetto_v1_0.pdf}
                    \item[] (Ultimo accesso: 14/04/2026)
                \end{itemize}
                \item Regolamento del Progetto Didattico
                \begin{itemize}
                    \item \url{https://www.math.unipd.it/~tullio/IS-1/2025/Dispense/PD1.pdf}
                    \item[] (Ultimo accesso: 24/04/2026)
                \end{itemize}
            \end{itemize}
            }
            \subsubsection{Riferimenti informativi} {
                \begin{itemize}
                \item \textit{Analisi dei Requisiti v1.0}
                \begin{itemize}
                    \item \url{https://team-three-technologies.github.io/Documenti/1%20-%20RTB/Documenti%20esterni/Analisi_dei_Requisiti_v1_0.pdf}
                    \item[] (Ultimo accesso: 24/04/2026)
                \end{itemize}
                \item Verbali interni
                \item Verbali esterni
                \item Capitolato C3
                \begin{itemize}
                    \item \url{https://www.math.unipd.it/~tullio/IS-1/2025/Progetto/C3.pdf}
                    \item[] (Ultimo accesso: 11/04/2026)
                \end{itemize}
                \item Slide del corso: T2 - Processi di ciclo di vita del \textit{software}
                \begin{itemize}
                    \item \url{https://www.math.unipd.it/~tullio/IS-1/2025/Dispense/T02.pdf}
                    \item[] (Ultimo accesso: 23/04/2026)
                \end{itemize}
                \item Slide del corso: T4 - Gestione di progetto
                \begin{itemize}
                    \item \url{https://www.math.unipd.it/~tullio/IS-1/2025/Dispense/T04.pdf}
                    \item[] (Ultimo accesso: 11/04/2026)
                \end{itemize}
                \item \textit{Glossario v0.5}
                \begin{itemize}
                    \item \url{https://team-three-technologies.github.io/Documenti/1%20-%20RTB/Documenti%20interni/Glossario_v0_5.pdf}
                    \item[] (Ultimo accesso: 24/04/2026)
                \end{itemize}
            \end{itemize}
            }
        }
    }
    
    \section{Installazione}
        \subsection{Requisiti di sistema}
        DIPReader è un'applicazione desktop autonoma che non richiede l'installazione di componenti aggiuntive, in quanto tutto il necessario per il funzionamento è già incluso nel pacchetto di distribuzione.\\

        L'applicazione è compatibile con i seguenti sistemi operativi:
        \begin{itemize}
            \item Windows: Windows 10 o versioni successive
            \item macOS: macOS 11 o versioni successive
            \item Linux: qualsiasi distribuzione moderna con supporto al formato AppImage
        \end{itemize}

        \subsection{Installazione}
            Per installare l'applicativo DIPReader la procedura cambia leggermente in base al sistema operativo in uso ma in nessun caso è richiesta l'installazione di componenti aggiuntive.
            L'applicativo è disponibile per il download alla seguente repository: \\
            \url{https://github.com/Team-Three-Technologies/MVP}

            (procedura con screen di come fare il download dalla repo)

            \begin{itemize}
                \item Windows: Il pacchetto è un file eseguibile .exe di tipo portable.
                    Per installare l'applicazione è sufficiente:
                    \begin{enumerate}
                        \item Posizionare il file scaricato nella cartella radice del pacchetto DIP
                        \item Avviare l'applicazione con un doppio clic
                    \end{enumerate}

                \item macOS: Il pacchetto per macOS è distribuito come archivio .zip. \\
                    Per installare l'applicazione è sufficiente:
                    \begin{enumerate}
                        \item Aprire il file .zip scaricato
                        \item Trascinare l'applicazione nella cartella radice del pacchetto DIP
                        \item Avviare l'applicazione con un doppio clic
                    \end{enumerate}

                \item Linux: Il pacchetto è distribuito nel formato AppImage, un formato universale compatibile con la maggior parte delle distribuzioni Linux. \\
                    Per avviare l'applicazione:
                    \begin{enumerate}
                        \item Scaricare il file .AppImage
                        \item Renderlo eseguibile: tasto destro sul file $\rightarrow$ Proprietà $\rightarrow$ spuntare "Consenti esecuzione come programma"
                        \item Spostare il file nella cartella radice del pacchetto DIP e avviarlo con un doppio clic
                    \end{enumerate}
            \end{itemize}

            In nessun caso sono richiesti permessi di amministratore.

            \subsection{Primo avvio}
            Al primo avvio DIPReader si apre direttamente sull'interfaccia principale, senza richiedere registrazione, creazione di account o configurazione. 
            Se la lista dei documenti risulta vuota, verificare che l'eseguibile sia posizionato correttamente nella cartella radice del pacchetto DIP (vedere sezione \ref{sec:apertura}).

            \begin{figure}[h!]
                \centering
                \setlength{\fboxrule}{0.5pt}
                \setlength{\fboxsep}{0pt}
                \fbox{
                    \includegraphics[width=0.5\textwidth]{screenshots_dipreader/dipreader.png}
                }
                \caption{Schermata principale}
                \label{fig:dipreader}
            \end{figure}


    \section{Guida all'utilizzo}
        \subsection{Apertura di un pacchetto DIP}
        \label{sec:apertura}
        DIPReader legge automaticamente il contenuto del pacchetto informativo all'avvio, senza necessità di configurazioni. Per caricare un DIP è sufficiente posizionare il file eseguibile dell'applicazione direttamente all'interno della cartella principale (o radice) del pacchetto DIP e avviarlo con un doppio clic.
        Il sistema ne riconosce automaticamente la struttura e popola la lista dei documenti. Nel caso di pacchetto non riconosciuto o con struttura non conforme, viene mostrato un apposito messaggio di errore.

        \subsection{Consultazione dei documenti}
            \subsubsection{Visualizzare la lista dei documenti}
            Al centro della finestra è presente la lista di tutti i documenti contenuti nel pacchetto DIP o, nel caso in cui siano applicati dei filtri (per maggiori dettagli sull'utilizzo dei filtri si rimanda alla sezione \ref{sec:filtri}), l'elenco dei soli documenti corrispondenti ai criteri di ricerca. Per ciascun 
            documento principale sono elencati, immediatamente al di sotto e con un leggero rientro, i relativi allegati.  
            
            \begin{figure}[h!]
                \centering
                \setlength{\fboxrule}{0.5pt}
                \setlength{\fboxsep}{0pt}
                \fbox{
                    \includegraphics[width=0.5\textwidth]{screenshots_dipreader/document-list.png}
                }
                \caption{Lista dei documenti}
                \label{fig:document-list}
            \end{figure}
            
            \subsubsection{Selezionare un documento e visualizzarne i dettagli}
            Per selezionare un documento è sufficiente fare clic sul nome del documento dalla lista. 
            Dopo aver selezionato un documento, nel pannello di destra compare la sezione "Dettagli Documento" con i seguenti campi: il nome del file, il formato (ad esempio PDF o XML), l'identificativo univoco del documento, e la data e ora di registrazione. Più in basso è presente la sezione "Allegati", che mostra il numero di allegati e li elenca ciascuno con il proprio percorso e formato.
            Per la consultazione degli allegati vedere la sezione~\ref{sec:allegati}.

            \begin{figure}[h!]
                \centering
                \setlength{\fboxrule}{0.5pt}
                \setlength{\fboxsep}{0pt}
                \fbox{
                    \includegraphics[width=0.5\textwidth]{screenshots_dipreader/document-details.png}
                }
                \caption{Dettagli di un documento}
                \label{fig:document-details}
            \end{figure}

            \subsubsection{Visualizzare l'anteprima di un documento}
            Per visualizzare il contenuto di un file, è sufficiente cliccare sull'icona a forma di occhio (\faEye) posta in corrispondenza del documento di interesse nell'elenco centrale. L'anteprima viene caricata in una schermata dedicata a tutto schermo, permettendo una lettura agevole senza la necessità di aprire programmi esterni. Per tornare alla lista di documenti è sufficiente cliccare l'apposito pulsante in alto a sinistra "Torna Indietro".
            Se l'anteprima non è disponibile, la sezione lo comunica tramite un apposito messaggio.

            \begin{figure}[h!]
                \centering
                \setlength{\fboxrule}{0.5pt}
                \setlength{\fboxsep}{0pt}
                \fbox{
                    \includegraphics[width=0.5\textwidth]{screenshots_dipreader/preview.png}
                }
                \caption{Anteprima di un documento}
                \label{fig:preview}
            \end{figure}

            \subsubsection{Esportare un documento}
            Per salvare una copia del documento sul proprio dispositivo, è sufficiente selezionare il documento desiderato dalla lista e fare clic sul pulsante "Esporta" presente nel pannello di destra. Verrà aperta una finestra di dialogo che permette di scegliere la cartella di destinazione.

        \subsection{Consultazione degli allegati}
            \subsubsection{Visualizzare gli allegati di un documento}
            \label{sec:allegati}
            Gli allegati sono accessibili direttamente dalla lista centrale dei documenti. Essi sono posizionati immediatamente al di sotto del rispettivo documento principale (padre), incolonnati e caratterizzati da un leggero rientro.
            
            \subsubsection{Selezionare un allegato e visualizzarne i dettagli}
            Per consultare i dettagli di uno specifico allegato, è sufficiente fare clic sul suo nome all'interno di questa sotto-lista. A seguito della selezione, il pannello di destra si aggiornerà mostrando la scheda "Dettagli Allegato", che riepiloga le informazioni specifiche del file selezionato quali: nome, formato e identificativo.

            \begin{figure}[h!]
                \centering
                \setlength{\fboxrule}{0.5pt}
                \setlength{\fboxsep}{0pt}
                \fbox{
                    \includegraphics[width=0.5\textwidth]{screenshots_dipreader/attachment.png}
                }
                \caption{Dettagli di un allegato}
                \label{fig:attachment}
            \end{figure}
            
            \subsubsection{Visualizzare l'anteprima di un allegato}
            Così come per i documenti principali, anche per gli allegati è possibile visualizzare l'anteprima. Per farlo, occorre cliccare sull'icona a forma di occhio (\faEye) situata alla destra del nominativo dell'allegato desiderato.

            \begin{figure}[h!]
                \centering
                \setlength{\fboxrule}{0.5pt}
                \setlength{\fboxsep}{0pt}
                \fbox{
                    \includegraphics[width=0.5\textwidth]{screenshots_dipreader/preview.png}
                }
                \caption{Anteprima di un documento}
            \end{figure}
            
            \subsubsection{Tornare al documento padre}
            Per tornare alla visualizzazione del documento principale è sufficiente fare clic sul suo nome nella lista centrale. Il pannello di destra si aggiorna ripristinando i metadati.

        \subsection{Ricerca e uso dei filtri}
            \label{sec:filtri}
                \subsubsection{Aggiungere un filtro}
                Nella parte superiore del pannello laterale sinistro è presente il pulsante "Aggiungi Filtro". Premendolo compare una nuova riga contenente un menu a tendina per selezionare il campo su cui effettuare la ricerca e un campo per inserire il valore desiderato.
                Quest'ultimo varia in base al campo selezionato: per i campi temporali è presente un selettore di data, per i campi categorici un menu a tendina con opzioni predefinite, per gli altri invece un campo libero.

            \begin{figure}[h!]
                \centering
                \setlength{\fboxrule}{0.5pt}
                \setlength{\fboxsep}{0pt}
                \fbox{
                    \includegraphics[width=0.5\textwidth]{screenshots_dipreader/filter1.png}
                }
                \caption{Aggiungere un filtro}
                \label{fig:filter1}
            \end{figure}

            \begin{figure}[h!]
                \centering
                \setlength{\fboxrule}{0.5pt}
                \setlength{\fboxsep}{0pt}
                \fbox{
                    \includegraphics[width=0.5\textwidth]{screenshots_dipreader/filter2.png}
                }
                \caption{Filtro con menu a tendina}
                \label{fig:filter2}
            \end{figure}

            \subsubsection{Utilizzare più filtri combinati}
            È possibile combinare più criteri di ricerca premendo nuovamente il pulsante "Aggiungi Filtro". I filtri aggiuntivi si applicano in modo cumulativo: la ricerca restituisce solo i documenti che soddisfano tutti i criteri inseriti contemporaneamente.

            \subsubsection{Rimuovere un filtro}
            \label{sec:removefilter}
            Per eliminare un filtro specifico è sufficiente fare clic sul pulsante a forma di "X" posto alla destra del relativo campo di testo.

            \begin{figure}[h!]
                \centering
                \setlength{\fboxrule}{0.5pt}
                \setlength{\fboxsep}{0pt}
                \fbox{
                    \includegraphics[width=0.5\textwidth]{screenshots_dipreader/filter3.png}
                }
                \caption{Pulsante per eliminare un criterio di ricerca attivo}
                \label{fig:filter3}
            \end{figure}

            \subsubsection{Avviare la ricerca}
            Una volta impostati i parametri desiderati, fare clic sul pulsante "Cerca" per avviare l'elaborazione. La lista si aggiorna mostrando solo i documenti che corrispondono esattamente ai criteri inseriti.

            \subsubsection{Ripristinare la lista completa}
            Per tornare a visualizzare l'intero contenuto del pacchetto DIP è sufficiente rimuovere tutti i filtri attivi. Per rimuovere un filtro si veda la sezione \ref{sec:removefilter}.
            
        \subsection{Visualizzare le informazioni sul pacchetto DIP}
        Il riquadro "Informazioni del pacchetto (DIP)", sempre visibile in basso a sinistra, riporta: UUID, data di creazione, numero totale di documenti contenuti e numero AIP.
        
        \begin{figure}[h!]
            \centering
            \setlength{\fboxrule}{0.5pt}
            \setlength{\fboxsep}{0pt}
            \fbox{
                \includegraphics[width=0.5\textwidth]{screenshots_dipreader/dip-details.png}
            }
            \caption{Dettagli del pacchetto DIP}
            \label{fig:dip-details}
        \end{figure}
\end{document}